\chapter{Text, graphemes, and width}
\label{ch:text}

\begin{aside}
A terminal puts characters in cells. Two facts make this harder
than it sounds: not every Unicode codepoint is one character, and
not every character is one cell wide. \maya{} treats text as a
grapheme stream and resolves visual width at the boundary between
the element tree and the canvas. This chapter explains the model.
\end{aside}

\section{Code units, codepoints, graphemes}

Three terms for three different things:

\begin{itemize}[leftmargin=*]
  \item A \textbf{code unit} is the smallest piece of an encoded
    string. For UTF-8, a code unit is a byte. For UTF-16, two
    bytes. \code{std::string::size()} returns code units.
  \item A \textbf{codepoint} is a Unicode scalar value
    (U+0000--U+10FFFF, minus the surrogate range). One codepoint
    may take one to four UTF-8 bytes.
  \item A \textbf{grapheme cluster} is what a human would call ``a
    character'' --- a base character plus any combining marks,
    joined emoji, and similar. One grapheme may contain several
    codepoints.
\end{itemize}

The terminal renders graphemes, but it does not actually know what
a grapheme is. It writes codepoints sequentially and lets the
font shaper combine them. If you put two codepoints in two cells
when they should have combined into one cell, the cursor and the
rest of the row drift by one.

\subsection*{An example}

The visible character \emph{é} can be encoded two ways:

\begin{itemize}[leftmargin=*]
  \item U+00E9 (precomposed): one codepoint, one cell.
  \item U+0065 U+0301 (e + combining acute): two codepoints, one
    cell.
\end{itemize}

\maya{} normalises text to a grapheme stream when measuring width
and emitting cells. \code{std::string} is opaque to all this; the
text module wraps it.

\begin{funfact}
Unicode has at least three different forms of \emph{é} when you
count the precomposed/decomposed variants. Normalisation forms
(NFC, NFD, NFKC, NFKD) reduce them to canonical representations.
\maya{} does not normalise inputs; it preserves the user's
bytes and walks them as graphemes.
\end{funfact}

\section{East Asian Width}

A second wrinkle: some codepoints take two cells. Chinese,
Japanese, Korean characters typically do; emoji do; many symbols
do. Unicode's \emph{East Asian Width} property classifies each
codepoint as Narrow, Wide, Ambiguous, Fullwidth, Halfwidth, or
Neutral.

\maya{} ships with a precomputed table in
\file{maya/include/maya/text/unicode\_width\_table.hpp} mapping
codepoint to cell width:

\begin{lstlisting}
int width_of(char32_t cp);
// 0 for combining marks, zero-width joiners, control chars
// 1 for narrow / neutral / halfwidth
// 2 for wide / fullwidth
\end{lstlisting}

The table is generated from the Unicode 15.x EastAsianWidth
data. Lookup is a binary search over the table; sub-microsecond.

\subsection*{Why the table is precomputed}

Computing width at runtime would mean walking the Unicode
character database, which is large. \maya{} ships the table as a
constexpr array; the lookup is fast and the binary stays
self-contained.

The trade-off: when Unicode releases a new version with new
codepoints, \maya{} needs a regenerate. The generator script
lives at \file{maya/tools/gen\_unicode\_width.py}.

\section{The grapheme walker}

\file{maya/include/maya/text/unicode\_width.hpp} exposes:

\begin{lstlisting}
class GraphemeWalker {
public:
    explicit GraphemeWalker(std::string_view text);

    struct Grapheme {
        std::string_view text;   // bytes for this grapheme
        int width;               // 0, 1, or 2
    };

    std::optional<Grapheme> next();
    void reset();
};
\end{lstlisting}

The walker steps through the input grapheme by grapheme. For
each grapheme it returns the byte span and the cell width.

\subsection*{Implementation sketch}

The walker is a small state machine driven by Unicode's grapheme
cluster boundaries:

\begin{itemize}[leftmargin=*]
  \item Skip combining marks (they belong to the previous
    grapheme).
  \item Handle regional indicators (flag emoji are pairs).
  \item Handle ZWJ sequences (joiner-joined emoji are one
    grapheme).
  \item Handle Hangul syllables (\emph{L V T} sequences are one).
\end{itemize}

The result is a stream of cleanly-bounded graphemes that the
canvas can lay out cell by cell.

\section{Measuring width}

The most common use:

\begin{lstlisting}
int display_width(std::string_view text) {
    GraphemeWalker w{text};
    int total = 0;
    while (auto g = w.next()) total += g->width;
    return total;
}
\end{lstlisting}

The layout phase calls this for each text node to get its
intrinsic width. Cached on the text node so successive frames do
not repeat the walk.

\subsection*{Ambiguous-width chars}

Some characters have ``ambiguous'' width: 1 cell in Western
terminals, 2 cells in CJK terminals. \maya{} defaults to 1
(Western), with a config switch for users on CJK terminals:

\begin{lstlisting}
maya::set_ambiguous_width(2);  // CJK terminals
\end{lstlisting}

The cached widths invalidate when the setting changes.

\section{Truncation and wrapping}

When a text node must fit a fixed width, \maya{} provides:

\begin{itemize}[leftmargin=*]
  \item \code{TextWrap::Truncate}: drop characters from the end,
    optionally with an ellipsis.
  \item \code{TextWrap::Wrap}: spill onto the next line at a
    word boundary.
  \item \code{TextWrap::WrapAnywhere}: spill at any grapheme
    boundary, including mid-word.
\end{itemize}

\subsection*{Truncation with ellipsis}

\begin{lstlisting}
std::string truncate(std::string_view text, int max_width) {
    GraphemeWalker w{text};
    std::string out;
    int used = 0;
    while (auto g = w.next()) {
        if (used + g->width > max_width - 1) {
            out += \"...\";
            return out;
        }
        out += g->text;
        used += g->width;
    }
    return out;
}
\end{lstlisting}

Reserve one cell for the ellipsis itself (well, three for the
ASCII ``...''; \maya{} can also use U+2026 for one cell).

\subsection*{Word wrapping}

Word wrap is harder. The algorithm:

\begin{enumerate}[leftmargin=*]
  \item Walk graphemes accumulating into a current line.
  \item When the line reaches \code{max\_width}, scan backward
    for the last whitespace and break there. The trailing word
    becomes the start of the next line.
  \item If a single word exceeds \code{max\_width}, break it
    at a grapheme boundary (degrade to WrapAnywhere).
\end{enumerate}

The result is a vector of lines. The canvas paints each line at
its computed $y$.

\section{Special characters}

Several Unicode categories deserve special handling.

\subsection*{Tabs}

A tab character (U+0009) advances the cursor to the next tab
stop, typically every 8 cells. \maya{} expands tabs to spaces
before measuring; the tab stops are configurable.

\subsection*{Newlines}

\code{\\n} (U+000A) starts a new line. The canvas treats it as a
forced break: the current row ends; the next grapheme goes to
column 0 of the next row.

\subsection*{Control characters}

Most control characters (U+0001--U+0008, U+000B--U+001F) are
zero-width and invisible. \maya{} either drops them or renders a
caret notation (\code{\textasciicircum A} for U+0001) depending
on the \code{ControlPolicy} config.

\subsection*{Variation selectors}

U+FE0E (text presentation) and U+FE0F (emoji presentation) modify
the preceding codepoint. \maya{} honours them: a base character
followed by U+FE0F becomes the emoji form.

\subsection*{Right-to-left text}

\maya{} does not perform full BiDi (bidirectional) reordering.
RTL text (Arabic, Hebrew) is rendered in logical order; the
terminal handles visual order based on its own BiDi
implementation. This means RTL text inside boxes may align
unexpectedly. Full BiDi is a planned feature; see the issue
tracker.

\begin{warning}
Mixing LTR and RTL in the same text node is a known weak point.
If you need this, render the RTL portion as its own text node
and accept the visual quirks.
\end{warning}

\section{Performance}

Width measurement is fast:

\begin{itemize}[leftmargin=*]
  \item Grapheme walking: roughly 10ns per grapheme.
  \item Width lookup: roughly 20ns per codepoint (binary search
    in the precomputed table).
  \item Per-line measurement: microseconds for 80-character
    lines.
\end{itemize}

For an 80x24 screen of text, total measurement is well under a
millisecond. \maya{} caches per-text-node width so re-measurement
across frames is free unless the text changes.

\subsection*{Cache invalidation}

The cache key is the text content plus the ambiguous-width
setting. If either changes, the cache is dropped.

\section{Pitfalls}

\begin{warning}
\textbf{Using std::string::size() for visual width.} It is the
byte count, not the cell count. ``hello'' is 5 bytes and 5
cells, but ``\verb|privet|'' (transliterating six Cyrillic
letters) is 12 bytes and 6 cells, and a two-character CJK string
like ``ni hao'' encoded in Chinese is 6 bytes and 4 cells. Use
\code{display\_width} for any
visual sizing.
\end{warning}

\begin{warning}
\textbf{Calling display\_width every frame.} For text in a tight
inner loop, the cost adds up. Pre-compute and cache, or trust
the layout phase to do it once.
\end{warning}

\begin{warning}
\textbf{Mixing widths in a single text node.} If a text node
contains both narrow and wide characters, the layout phase
handles it. But if you try to truncate at a fixed
character count instead of cell count, you will get visual
misalignment. Always truncate by cell count.
\end{warning}

\begin{warning}
\textbf{Trailing combining marks.} A grapheme can extend past
the visual edge of a fixed-width region. The walker handles
this, but if you slice raw bytes you may split a grapheme. Always
slice at grapheme boundaries.
\end{warning}

\section*{Workshop}

\begin{enumerate}[leftmargin=*]
  \item Write a function that takes a \code{std::string\_view}
    and returns its display width using \code{GraphemeWalker}.
    Test it on plain ASCII, on emoji, and on CJK text.
  \item Write a truncator that respects grapheme boundaries.
    Test it on a string ending in a multi-codepoint emoji.
  \item Read
    \file{maya/include/maya/text/unicode\_width\_table.hpp}.
    Note the size of the table and the structure of its
    entries.
\end{enumerate}

\begin{recap}
Text is a stream of bytes, codepoints, and graphemes; the cell
grid wants graphemes. \maya{} ships a precomputed width table
and a grapheme walker that converts. Truncation and wrapping
respect grapheme boundaries. Tabs, newlines, control characters,
variation selectors, and ambiguous-width chars have explicit
policies.
\end{recap}
